\section{Discussion}

Our experiments validate the fingerprinting hypothesis while revealing important scope limitations.

\subsection{Key Findings}

\textbf{Attribution Methods Are Radically Different.} The F-ratio of 1423.55 far exceeds what we anticipated. This indicates that method choice is not a minor implementation detail---different methods measure fundamentally different aspects of training data influence.

\textbf{Fingerprints Enable Method Selection.} The stability of mode profiles ($\alpha > 0.96$) means fingerprints are reliable signatures. This opens a path toward principled method selection: memorization detection $\rightarrow$ TRAK; curvature analysis $\rightarrow$ Kronfluence.

\textbf{Architecture Dependence Bounds Generalization.} The ConvNeXt inversion is our most significant negative result. Fingerprints transfer within architectural families but invert across fundamentally different architectures.

\subsection{Limitations}

\textbf{L1: Architecture-Dependent Transfer.} Mode profiles do not transfer universally. Results apply to within-family comparisons; cross-family requires calibration.

\textbf{L2: CPU-Only Execution.} Due to CUDA driver incompatibility, experiments ran on CPU with reduced scale. Effect sizes far exceed thresholds, suggesting conclusions are robust.

\textbf{L3: Vision Proxy for LLM Hypothesis.} The original motivation involved LLM attribution, but validation used vision models. The mechanism should transfer theoretically.

\subsection{Broader Impact}

Our fingerprinting framework can improve data attribution practice by enabling informed method selection for fairness audits, data valuation, and model debugging.
